\chapter{The identity component \(\Pic^0_{C/k}\) as an abelian variety}
\label{chap:Picard_Pic0AbelianVariety}
% archon:covers AlgebraicJacobian/Picard/Pic0AbelianVariety.lean

% Chapter-local macro: \tu renders an upright (roman) glyph in math mode,
% used here only for the tangent-space symbol T_0 (Kleiman's \tu = \textup).
\providecommand{\tu}{\textup}

% STRATEGY NOTE (iter-192). This chapter packages strategy phase
% A.3.iii--vi (tangent space, smoothness, properness, geometric
% irreducibility) and their A.3.vii assembly into the abelian-variety
% statement \(\Pic^0_{C/k}\) is an abelian variety, the load-bearing
% gate of Route~A for \texttt{thm:nonempty\_jacobianWitness}.
%
% The chapter consumes the abstract identity-component substrate set
% up in \cref{chap:Picard_IdentityComponent} (\(\Pic^0_{C/k}\) as a
% clopen subgroup scheme of finite type, geometrically irreducible,
% formation commuting with base change) and adds the curve-specific
% input: the tangent-space isomorphism
% \(\tu T_0 \Pic^0_{C/k} \cong H^1(C, \mathcal{O}_C)\), the resulting
% smoothness of dimension \(g\) at the origin (hence everywhere by
% translation), properness inherited from the projectivity of
% \(\Pic^0_{C/k}\) for \(C/k\) geometrically normal, and the assembly
% into an abelian variety in the sense of Milne~\S I.1.
%
% NOTE (planner directive iter-192). The Lean target file
% \texttt{AlgebraicJacobian/Picard/Pic0AbelianVariety.lean} does not
% yet exist; this chapter is the authoritative informal specification,
% and the Lean skeleton is owed in a follow-up iteration.
%
% NOTE (run-0008). After the FGA picSharp rewire, Pic0Scheme C is a REAL
% definition (GroupScheme.IdentityComponent (PicScheme C)), and every Lean
% statement of this chapter that mentions Pic0Scheme C now carries the
% instance hypotheses [HasPicScheme C] [PicScheme.PicSchemeLocallyOfFiniteType C]
% (supplied under a k-rational point via Scheme.HasRationalPoint; see
% chap:Picard_FGAPicRepresentability, def:has_rational_point /
% def:pic_scheme_lft). Mathematically: the chapter's statements hold for a
% pointed curve, per the route-(b) convention of the FGA chapter. The
% tangent-space isomorphism can now attach to the genuine bijection
% (T -> Pic_{C/k}) ≃ Pic(C ×_k T)/π_T^* Pic(T) extracted from
% PicScheme.representable at T = Spec k[ε].

\section{Tangent space at the identity}
\label{sec:pic0_tangent_space}

For a \(k\)-group scheme \(G\) locally of finite type, the tangent
space at the identity \(\tu T_0 G\) is the \(k\)-vector space of
\(k\)-derivations of its local ring at \(e\) into \(k\), or
equivalently (and functorially) the set
\(G(k[\varepsilon]/(\varepsilon^2))_e\) of \(k_\varepsilon\)-points
whose underlying \(k\)-point is \(e\). For the Picard scheme of a
geometrically integral projective \(k\)-scheme, this tangent space is
identified with the first sheaf cohomology of the structure sheaf via
the truncated exponential sequence and the étale-sheaf comparison
theorem; on a smooth proper curve, this is the genus-many-dimensional
\(k\)-vector space \(H^1(C, \mathcal{O}_C)\).

We first record the local-algebra core of the "tangent space via dual
numbers" identification: for a local \(k\)-algebra \(R\) with maximal
ideal \(\mathfrak m\) and residue field \(\kappa\), and with the
\(k\)-rational-point hypothesis that the structure map
\(k \to \kappa\) is an isomorphism, the \(k\)-derivations of \(R\)
into \(\kappa\), the local \(k\)-algebra maps
\(R \to k[\varepsilon]/(\varepsilon^2)\), and the linear functionals
on the cotangent space \(\mathfrak m/\mathfrak m^2\) all agree.
Applied to the local ring of \(\Pic_{C/k}\) at the identity, this is
the identification \(\tu T_0 \Pic_{C/k} \cong
\Hom(\mathfrak m_0/\mathfrak m_0^2, k)\) used in the main theorem.

\begin{lemma}[Derivations descend to the cotangent space]\leanok
  \label{lem:tangent_derivation_descends}
  \dcref{custom}
  \lean{AlgebraicGeometry.derivationCotangentDual}
  Let \(k\) be a field and \(R\) a local \(k\)-algebra with maximal
  ideal \(\mathfrak m\) and residue field \(\kappa\). Every
  \(k\)-derivation \(d : R \to \kappa\) vanishes on
  \(\mathfrak m^2\), and its restriction to \(\mathfrak m\) is
  \(R\)-linear; consequently \(d\) induces an \(R\)-linear map
  \[
    \bar d : \mathfrak m/\mathfrak m^2 \longrightarrow \kappa .
  \]
\end{lemma}

\begin{proof}\leanok
  Every element of \(\mathfrak m\) acts by zero on the \(R\)-module
  \(\kappa = R/\mathfrak m\), since the action factors through the
  residue class. For \(x, y \in \mathfrak m\) the Leibniz rule gives
  \(d(xy) = x \cdot d(y) + y \cdot d(x) = 0\), so \(d\) vanishes on
  \(\mathfrak m^2\). For \(r \in R\) and \(x \in \mathfrak m\),
  \(d(rx) = r \cdot d(x) + x \cdot d(r) = r \cdot d(x)\), so
  \(d|_{\mathfrak m}\) is \(R\)-linear; it therefore descends to the
  quotient \(\mathfrak m/\mathfrak m^2\).
\end{proof}

\begin{lemma}[Derivations are the cotangent dual]\leanok
  \label{lem:tangent_derivation_cotangent_dual}
  \dcref{custom}
  \lean{AlgebraicGeometry.derivationEquivCotangentDual}
  \uses{lem:tangent_derivation_descends}
  Let \(k\) be a field and \(R\) a local \(k\)-algebra whose
  structure map \(k \to \kappa\) to the residue field is an
  isomorphism. Then descent \(d \mapsto \bar d\)
  (\cref{lem:tangent_derivation_descends}) is an isomorphism of
  abelian groups
  \[
    \operatorname{Der}_k(R, \kappa)
    \;\cong\;
    \Hom_R(\mathfrak m/\mathfrak m^2, \kappa).
  \]
\end{lemma}

\begin{proof}\leanok
  Write \(\sigma : \kappa \to k\) for the inverse of the structure
  map. Each \(r \in R\) splits canonically as
  \(r = c_r + x_r\) with constant part
  \(c_r := \sigma(\bar r) \cdot 1 \in k \cdot 1\) and
  \(x_r := r - c_r \in \mathfrak m\) (its residue class vanishes by
  construction).

  \emph{Injectivity.} A \(k\)-derivation kills the constants
  \(k \cdot 1\), so \(d(r) = d(x_r)\) for every \(r\): the derivation
  is determined by its restriction to \(\mathfrak m\), i.e.\ by
  \(\bar d\).

  \emph{Surjectivity.} Given an \(R\)-linear
  \(\varphi : \mathfrak m/\mathfrak m^2 \to \kappa\), set
  \(d(r) := \varphi(x_r \bmod \mathfrak m^2)\). Additivity and
  \(k\)-linearity follow since \(r \mapsto c_r\) is \(k\)-linear. For
  the Leibniz rule, expand
  \[
    x_{rs} \;=\; c_r\, x_s + c_s\, x_r + x_r x_s ,
  \]
  where \(x_r x_s \in \mathfrak m^2\) dies in the cotangent space;
  since \(r\) and \(c_r\) have the same residue class, they act
  identically on \(\kappa\), whence
  \(d(rs) = c_r \varphi(\bar x_s) + c_s \varphi(\bar x_r)
    = r \cdot d(s) + s \cdot d(r)\).
  On elements of \(\mathfrak m\) the constant part vanishes, so the
  two constructions are mutually inverse.
\end{proof}

\begin{lemma}[Dual-number points are derivations]\leanok
  \label{lem:tangent_dual_number_derivation}
  \dcref{custom}
  \lean{AlgebraicGeometry.localDualNumberHomEquivDerivation}
  Let \(k\) be a field and \(R\) a local \(k\)-algebra whose
  structure map \(k \to \kappa\) is an isomorphism. Sending a
  \(k\)-algebra homomorphism \(f : R \to k[\varepsilon]/(\varepsilon^2)\)
  to the \(\varepsilon\)-coefficient of its values is a bijection
  \[
    \{\, f : R \to k_\varepsilon \ \text{local $k$-algebra maps} \,\}
    \;\cong\;
    \operatorname{Der}_k(R, \kappa),
  \]
  where locality means \(f(\mathfrak m) \subseteq (\varepsilon)\),
  the maximal ideal of the local ring \(k_\varepsilon\).
\end{lemma}

\begin{proof}\leanok
  Write \(f(r) = a(r) + b(r)\varepsilon\). Splitting
  \(r = c_r + x_r\) into constant part and maximal-ideal part as in
  \cref{lem:tangent_derivation_cotangent_dual}, locality and
  \(k\)-linearity force \(a(r) = a(c_r) = \sigma(\bar r)\): the
  constant component of \(f\) is the residue map read through
  \(\kappa \cong k\). Multiplicativity of \(f\) and
  \(\varepsilon^2 = 0\) give
  \(b(rs) = a(r)\,b(s) + a(s)\,b(r)\), which — transported through
  \(\kappa \cong k\) — is exactly the Leibniz rule for the map
  \(D := \bar b : R \to \kappa\); \(k\)-linearity of \(D\) is
  inherited from \(f\). Conversely, a derivation
  \(D : R \to \kappa\) assembles with the residue map into
  \(f(r) := \sigma(\bar r) + \sigma(D(r))\,\varepsilon\), a local
  \(k\)-algebra homomorphism by the same computation. The two
  constructions are mutually inverse: the constant component of any
  local \(f\) is forced as above, and the \(\varepsilon\)-component
  is preserved by both directions.
\end{proof}

\begin{corollary}[The tangent space is the cotangent dual]\leanok
  \label{cor:tangent_dual_number_cotangent_dual}
  \dcref{custom}
  \lean{AlgebraicGeometry.localDualNumberHomEquivCotangentDual}
  \uses{lem:tangent_derivation_cotangent_dual,
        lem:tangent_dual_number_derivation}
  Let \(k\) be a field and \(R\) a local \(k\)-algebra whose
  structure map \(k \to \kappa\) is an isomorphism. Then the set of
  local \(k\)-algebra maps \(R \to k_\varepsilon\) — the dual-number
  points of \(\operatorname{Spec} R\) lifting the closed point — is
  in bijection with
  \(\Hom_R(\mathfrak m/\mathfrak m^2, \kappa)\), the dual of the
  cotangent space.
\end{corollary}

\begin{proof}\leanok
  Compose the bijections of
  \cref{lem:tangent_dual_number_derivation} and
  \cref{lem:tangent_derivation_cotangent_dual}.
\end{proof}

\begin{lemma}[The cotangent dual is $\kappa$-linear]\leanok
  \label{lem:tangent_cotangent_dual_kappa_linear}
  \dcref{custom}
  \lean{AlgebraicGeometry.cotangentDualExtendScalars}
  Let \(R\) be a local ring with maximal ideal \(\mathfrak m\) and
  residue field \(\kappa\). Every \(R\)-linear map
  \(\mathfrak m/\mathfrak m^2 \to \kappa\) is \(\kappa\)-linear for
  the natural \(\kappa\)-vector-space structures on both sides:
  \[
    \Hom_R(\mathfrak m/\mathfrak m^2, \kappa)
    \;=\;
    \Hom_\kappa(\mathfrak m/\mathfrak m^2, \kappa)
    \;=\; (\mathfrak m/\mathfrak m^2)^{\vee}.
  \]
\end{lemma}

\begin{proof}\leanok
  Both \(\mathfrak m/\mathfrak m^2\) and \(\kappa\) are modules over
  \(\kappa = R/\mathfrak m\), and their \(R\)-actions factor through
  the surjection \(R \twoheadrightarrow \kappa\). Hence for an
  \(R\)-linear \(\varphi\) and a scalar \(c = \bar r \in \kappa\),
  \(\varphi(c \cdot x) = \varphi(r \cdot x) = r \cdot \varphi(x)
    = c \cdot \varphi(x)\).
\end{proof}

\begin{corollary}[Tangent space as the dual vector space of
  $\mathfrak m/\mathfrak m^2$]\leanok
  \label{cor:tangent_dual_number_cotangent_space_dual}
  \dcref{custom}
  \lean{AlgebraicGeometry.localDualNumberHomEquivCotangentSpaceDual}
  \uses{cor:tangent_dual_number_cotangent_dual,
        lem:tangent_cotangent_dual_kappa_linear}
  In the situation of \cref{cor:tangent_dual_number_cotangent_dual},
  the dual-number points of \(\operatorname{Spec} R\) lifting the
  closed point are in bijection with the \(\kappa\)-linear dual
  vector space \((\mathfrak m/\mathfrak m^2)^{\vee}\). In particular
  their \(\kappa\)-dimension equals
  \(\dim_\kappa \mathfrak m/\mathfrak m^2\) when the latter is
  finite.
\end{corollary}

\begin{proof}\leanok
  Compose \cref{cor:tangent_dual_number_cotangent_dual} with the
  identification of \cref{lem:tangent_cotangent_dual_kappa_linear};
  the dimension statement is the equality
  \(\dim_\kappa V^{\vee} = \dim_\kappa V\) for finite-dimensional
  \(V\).
\end{proof}

The local-algebra bijections above consume local homomorphisms out of
the stalk \(\mathcal O_{X,x}\). The next three statements supply the
geometric input: morphisms \(\operatorname{Spec} k_\varepsilon \to X\)
landing at a fixed point \(x\) are exactly such stalk data.

\begin{lemma}[Local maps to the dual numbers]\leanok
  \label{lem:tangent_dual_number_local_iff}
  \dcref{custom}
  \lean{AlgebraicGeometry.isLocalHom_dualNumber_iff}
  Let \(k\) be a field, \(R\) a local ring, and
  \(f : R \to k_\varepsilon = k[\varepsilon]/(\varepsilon^2)\) a ring
  homomorphism. Then \(f\) is a local homomorphism if and only if the
  constant component of \(f\) kills the maximal ideal:
  \[
    f(\mathfrak m) \subseteq (\varepsilon)
    \quad\Longleftrightarrow\quad
    \forall\, x \in \mathfrak m,\ \operatorname{fst}(f(x)) = 0 .
  \]
\end{lemma}

\begin{proof}\leanok
  The dual numbers \(k_\varepsilon\) form a local ring whose maximal
  ideal is \((\varepsilon)\), the kernel of the constant-component
  projection \(\operatorname{fst}\); an element of
  \(k_\varepsilon\) is a unit iff its constant component is a nonzero
  element of the field \(k\). If \(f\) is local and
  \(x \in \mathfrak m\), then \(x\) is a nonunit, so \(f(x)\) is a
  nonunit, i.e.\ \(\operatorname{fst}(f(x)) = 0\). Conversely, if the
  constant component kills \(\mathfrak m\) and \(f(a)\) is a unit,
  then \(\operatorname{fst}(f(a)) \neq 0\), so \(a \notin \mathfrak m\),
  i.e.\ \(a\) is a unit.
\end{proof}

\begin{lemma}[$\operatorname{Spec}$ of a local ring at a fixed point]\leanok
  \label{lem:tangent_spec_local_ring_at}
  \dcref{custom}
  \lean{AlgebraicGeometry.specToEquivOfLocalRingAt}
  Let \(R\) be a local ring, \(X\) a scheme, and \(x \in X\) a fixed
  point. Morphisms \(\operatorname{Spec} R \to X\) carrying the closed
  point of \(\operatorname{Spec} R\) to \(x\) are in bijection with the
  local ring homomorphisms \(\mathcal O_{X,x} \to R\):
  \[
    \{\, f : \operatorname{Spec} R \to X \mid f(\mathfrak m_R) = x \,\}
    \;\cong\;
    \{\, \varphi : \mathcal O_{X,x} \to R \ \text{local} \,\}.
  \]
\end{lemma}

\begin{proof}\leanok
  This is the fiber at \(x\) of the standard description of morphisms
  out of the spectrum of a local ring (Stacks 01J6; in Mathlib,
  \texttt{SpecToEquivOfLocalRing}):
  \(\operatorname{Hom}(\operatorname{Spec} R, X)\) is the set of pairs
  \((y, \varphi)\) of a point \(y \in X\) and a local homomorphism
  \(\varphi : \mathcal O_{X,y} \to R\). Forward, a morphism \(f\) with
  \(f(\mathfrak m_R) = x\) induces on stalks a local homomorphism
  \(\mathcal O_{X,f(\mathfrak m_R)} \to \mathcal O_{\operatorname{Spec} R,\mathfrak m_R}
    = R\), transported to \(\mathcal O_{X,x}\) along
  \(f(\mathfrak m_R) = x\). Backward, a local \(\varphi\) yields
  \(\operatorname{Spec}(\varphi)\) composed with the canonical map
  \(\operatorname{Spec} \mathcal O_{X,x} \to X\), which carries the
  closed point to \(x\) precisely because \(\varphi\) is local. The
  two constructions are mutually inverse by the same computation as
  for the unfibered bijection, the transport along \(f(\mathfrak m_R) = x\)
  cancelling on both sides.
\end{proof}

\begin{corollary}[Dual-number points at a fixed point]\leanok
  \label{cor:tangent_spec_dual_number_at}
  \dcref{custom}
  \lean{AlgebraicGeometry.specDualNumberAtEquiv}
  \uses{lem:tangent_dual_number_local_iff,
        lem:tangent_spec_local_ring_at}
  Let \(k\) be a field, \(X\) a scheme, and \(x \in X\). The
  dual-number points of \(X\) at \(x\) — morphisms
  \(\operatorname{Spec} k_\varepsilon \to X\) carrying the closed point
  to \(x\) — are in bijection with the ring homomorphisms
  \(\mathcal O_{X,x} \to k_\varepsilon\) whose constant component kills
  the maximal ideal. This is exactly the interface consumed by
  \cref{cor:tangent_dual_number_cotangent_space_dual}, once the
  \(k\)-algebra structure on the stalk at a \(k\)-rational point of a
  \(k\)-scheme is threaded through.
\end{corollary}

\begin{proof}\leanok
  Instantiate \cref{lem:tangent_spec_local_ring_at} at
  \(R = k_\varepsilon\) (a local ring, its maximal ideal being
  \((\varepsilon)\)) and rewrite the locality condition by
  \cref{lem:tangent_dual_number_local_iff}.
\end{proof}

The stalk data above is so far only ring-theoretic. The structure
morphism of a \(k\)-scheme endows each stalk with a canonical
\(k\)-algebra structure, and morphisms \emph{over} \(\operatorname{Spec} k\)
are exactly those whose stalk data is \(k\)-linear. This is the missing
thread between \cref{cor:tangent_spec_dual_number_at} and
\cref{cor:tangent_dual_number_cotangent_space_dual}.

\begin{definition}[Structure homomorphism into a stalk]\leanok
  \label{def:tangent_stalk_structure_hom}
  \dcref{custom}
  \lean{AlgebraicGeometry.stalkStructureHom}
  Let \(k\) be a field, \(X\) a scheme, \(f : X \to \operatorname{Spec} k\)
  a morphism, and \(x \in X\). The \emph{structure homomorphism}
  \(\sigma_x : k \to \mathcal O_{X,x}\) is the composite
  \[
    k \;\cong\; \Gamma(\operatorname{Spec} k, \mathcal O)
      \;\longrightarrow\; \mathcal O_{\operatorname{Spec} k,\, f(x)}
      \;\longrightarrow\; \mathcal O_{X,x}
  \]
  of the global-sections identification, the germ map at \(f(x)\), and
  the stalk map of \(f\) at \(x\).
\end{definition}

\begin{definition}[The stalk as a $k$-algebra]\leanok
  \label{def:tangent_stalk_algebra}
  \dcref{custom}
  \lean{AlgebraicGeometry.stalkAlgebra}
  \uses{def:tangent_stalk_structure_hom}
  In the situation of \cref{def:tangent_stalk_structure_hom}, the ring
  homomorphism \(\sigma_x\) makes the stalk \(\mathcal O_{X,x}\) a
  \(k\)-algebra, with \(\operatorname{algebraMap} = \sigma_x\).
\end{definition}

\begin{lemma}[Canonical map from the spectrum of a stalk, over $\operatorname{Spec} k$]\leanok
  \label{lem:tangent_from_spec_stalk_over}
  \dcref{custom}
  \lean{AlgebraicGeometry.fromSpecStalk_comp_eq}
  \uses{def:tangent_stalk_structure_hom}
  In the situation of \cref{def:tangent_stalk_structure_hom}, the
  canonical morphism
  \(c_x : \operatorname{Spec} \mathcal O_{X,x} \to X\) followed by \(f\)
  is the morphism of affine schemes induced by the structure
  homomorphism:
  \[
    f \circ c_x \;=\; \operatorname{Spec}(\sigma_x).
  \]
\end{lemma}

\begin{proof}\leanok
  By naturality of the canonical map out of the spectrum of a stalk,
  \(f \circ c_x\) factors as
  \(c_{f(x)} \circ \operatorname{Spec}(f^\sharp_x)\), where
  \(f^\sharp_x : \mathcal O_{\operatorname{Spec} k, f(x)} \to \mathcal O_{X,x}\)
  is the stalk map of \(f\) and
  \(c_{f(x)} : \operatorname{Spec} \mathcal O_{\operatorname{Spec} k, f(x)}
  \to \operatorname{Spec} k\) is the canonical map for the affine scheme
  \(\operatorname{Spec} k\). The latter is
  \(\operatorname{Spec}\) applied to the composite of the
  global-sections identification
  \(k \cong \Gamma(\operatorname{Spec} k, \mathcal O)\) with the germ
  map at \(f(x)\). By contravariant functoriality of
  \(\operatorname{Spec}\), the composite is
  \(\operatorname{Spec}\) of
  \(f^\sharp_x \circ (\text{germ at } f(x)) \circ (\text{identification})
  = \sigma_x\).
\end{proof}

\begin{lemma}[Over-triangles are structure-compatibilities]\leanok
  \label{lem:tangent_over_triangle_iff}
  \dcref{custom}
  \lean{AlgebraicGeometry.comp_eq_spec_iff}
  \uses{def:tangent_stalk_structure_hom}
  Let \(k\) be a field, \(f : X \to \operatorname{Spec} k\) a morphism
  of schemes, \(S\) a local ring, \(g : \operatorname{Spec} S \to X\),
  and \(\psi : k \to S\) a ring homomorphism. Write
  \(x = g(\mathfrak m_S)\) for the image of the closed point and
  \(\varphi_g : \mathcal O_{X,x} \to S\) for the local homomorphism
  induced by \(g\) on stalks
  (\cref{lem:tangent_spec_local_ring_at}). Then
  \[
    f \circ g = \operatorname{Spec}(\psi)
    \quad\Longleftrightarrow\quad
    \varphi_g \circ \sigma_x = \psi .
  \]
\end{lemma}

\begin{proof}\leanok
  \uses{lem:tangent_from_spec_stalk_over, lem:tangent_spec_local_ring_at}
  By \cref{lem:tangent_spec_local_ring_at} the morphism \(g\) factors
  as \(g = c_x \circ \operatorname{Spec}(\varphi_g)\) with
  \(c_x : \operatorname{Spec} \mathcal O_{X,x} \to X\) the canonical
  map. Hence, using \cref{lem:tangent_from_spec_stalk_over} and
  contravariant functoriality,
  \[
    f \circ g
      = (f \circ c_x) \circ \operatorname{Spec}(\varphi_g)
      = \operatorname{Spec}(\sigma_x) \circ \operatorname{Spec}(\varphi_g)
      = \operatorname{Spec}(\varphi_g \circ \sigma_x).
  \]
  Since \(\operatorname{Spec}\) is faithful on ring homomorphisms (it
  is fully faithful onto affine schemes), the equality
  \(f \circ g = \operatorname{Spec}(\psi)\) holds if and only if
  \(\varphi_g \circ \sigma_x = \psi\).
\end{proof}

\begin{corollary}[Over-$\operatorname{Spec} k$ dual-number points are local $k$-algebra maps]\leanok
  \label{cor:tangent_over_dual_number_alghom}
  \dcref{custom}
  \lean{AlgebraicGeometry.overDualNumberAtEquivAlgHom}
  \uses{def:tangent_stalk_algebra,
        cor:tangent_spec_dual_number_at,
        lem:tangent_over_triangle_iff}
  Let \(k\) be a field, \(X\) a scheme over \(\operatorname{Spec} k\),
  and \(x \in X\). The dual-number points of \(X\) at \(x\) \emph{over}
  \(\operatorname{Spec} k\) — morphisms
  \(g : \operatorname{Spec} k_\varepsilon \to X\) over
  \(\operatorname{Spec} k\) carrying the closed point to \(x\) — are in
  bijection with the \(k\)-\emph{algebra} homomorphisms
  \(\mathcal O_{X,x} \to k_\varepsilon\) whose constant component kills
  the maximal ideal. (A plain ring homomorphism
  \(\mathcal O_{X,x} \to k_\varepsilon\) need not be \(k\)-linear; the
  over-structure is what enforces \(k\)-linearity.)
\end{corollary}

\begin{proof}\leanok
  Under the bijection of \cref{cor:tangent_spec_dual_number_at},
  a morphism \(g\) with \(g(\mathfrak m) = x\) corresponds to the ring
  homomorphism \(\varphi_g : \mathcal O_{X,x} \to k_\varepsilon\), and
  by \cref{lem:tangent_over_triangle_iff} (applied with
  \(S = k_\varepsilon\) and \(\psi\) the canonical inclusion
  \(k \to k_\varepsilon\)) the over-triangle condition
  \(f \circ g = \operatorname{Spec}(k \to k_\varepsilon)\) holds if and
  only if \(\varphi_g \circ \sigma_x\) is the canonical inclusion,
  i.e.\ precisely when \(\varphi_g\) is a homomorphism of
  \(k\)-algebras for the structure of
  \cref{def:tangent_stalk_algebra}. The side condition on the constant
  component is untouched by this refinement.
\end{proof}

\begin{theorem}[Tangent space of a $k$-scheme at a rational point]\leanok
  \label{thm:tangent_over_dual_number_cotangent_dual}
  \dcref{custom}
  \lean{AlgebraicGeometry.overDualNumberAtEquivCotangentSpaceDual}
  \uses{cor:tangent_over_dual_number_alghom,
        cor:tangent_dual_number_cotangent_space_dual}
  Let \(k\) be a field, \(X\) a scheme over \(\operatorname{Spec} k\),
  and \(x \in X\) a \(k\)-rational point (the structure map
  \(k \to \kappa(x)\) to the residue field is bijective). Then the
  dual-number points of \(X\) at \(x\) over \(\operatorname{Spec} k\)
  are in bijection with the \(\kappa(x)\)-linear dual of the cotangent
  space:
  \[
    \{\, g : \operatorname{Spec} k_\varepsilon \to X
        \text{ over } \operatorname{Spec} k
        \mid g(\mathfrak m) = x \,\}
    \;\cong\;
    \operatorname{Hom}_{\kappa(x)}\!\bigl(\mathfrak m_x/\mathfrak m_x^2,\,
      \kappa(x)\bigr).
  \]
\end{theorem}

\begin{proof}\leanok
  Compose the bijection of \cref{cor:tangent_over_dual_number_alghom}
  with the local-algebra identification of
  \cref{cor:tangent_dual_number_cotangent_space_dual}, applied to the
  local \(k\)-algebra \(R = \mathcal O_{X,x}\)
  (\cref{def:tangent_stalk_algebra}), whose residue field is \(k\) by
  the rationality hypothesis.
\end{proof}

\begin{lemma}[A section retracts the structure homomorphism]\leanok
  \label{lem:tangent_section_retraction}
  \dcref{custom}
  \lean{AlgebraicGeometry.stalkStructureHom_comp_stalkClosedPointTo}
  \uses{def:tangent_stalk_structure_hom}
  Let \(k\) be a field, \(f : X \to \operatorname{Spec} k\) a morphism
  of schemes, and \(e : \operatorname{Spec} k \to X\) a section of
  \(f\) (that is, \(f \circ e = \operatorname{id}\)). Write
  \(x = e(*)\) for the image of the unique point of
  \(\operatorname{Spec} k\), \(\sigma_x : k \to \mathcal O_{X,x}\) for
  the structure homomorphism
  (\cref{def:tangent_stalk_structure_hom}), and
  \(\varphi_e : \mathcal O_{X,x} \to k\) for the local homomorphism
  induced by \(e\) on stalks. Then
  \(\varphi_e \circ \sigma_x = \operatorname{id}_k\).
\end{lemma}

\begin{proof}\leanok
  \uses{lem:tangent_over_triangle_iff}
  Apply \cref{lem:tangent_over_triangle_iff} with \(S = k\),
  \(g = e\) and \(\psi = \operatorname{id}_k\): the over-triangle
  \(f \circ e = \operatorname{id} =
  \operatorname{Spec}(\operatorname{id}_k)\) holds by hypothesis, and
  it is equivalent to \(\varphi_e \circ \sigma_x =
  \operatorname{id}_k\).
\end{proof}

\begin{lemma}[The image of a section is a rational point]\leanok
  \label{lem:tangent_section_rational}
  \dcref{custom}
  \lean{AlgebraicGeometry.bijective_algebraMap_residueField_of_section}
  \uses{def:tangent_stalk_algebra}
  Let \(k\) be a field, \(X\) a scheme over \(\operatorname{Spec} k\),
  \(e : \operatorname{Spec} k \to X\) a section of the structure
  morphism, and \(x = e(*)\). Then \(x\) is a \(k\)-rational point:
  the structure map \(k \to \kappa(x)\) to the residue field is
  bijective.
\end{lemma}

\begin{proof}\leanok
  \uses{lem:tangent_section_retraction}
  Injectivity is automatic: a ring homomorphism out of a field into a
  nonzero ring is injective. For surjectivity, the retraction
  \(\varphi_e : \mathcal O_{X,x} \to k\) of
  \cref{lem:tangent_section_retraction} is a local homomorphism into
  a field, so it kills the maximal ideal \(\mathfrak m_x\) and
  descends to a field homomorphism
  \(\bar\varphi : \kappa(x) \to k\) with
  \(\bar\varphi \circ \iota = \operatorname{id}_k\), where
  \(\iota : k \to \kappa(x)\) is the structure map. Given
  \(a \in \kappa(x)\), the element
  \(a - \iota(\bar\varphi(a))\) lies in the kernel of
  \(\bar\varphi\); but \(\bar\varphi\) is a homomorphism of fields,
  hence injective, so \(a = \iota(\bar\varphi(a))\). Thus \(\iota\)
  is surjective.
\end{proof}

\begin{corollary}[Tangent space at the image of a section]\leanok
  \label{cor:tangent_section_cotangent_dual}
  \dcref{custom}
  \lean{AlgebraicGeometry.overDualNumberSectionEquivCotangentSpaceDual}
  \uses{lem:tangent_section_rational}
  Let \(k\) be a field, \(X\) a scheme over \(\operatorname{Spec} k\),
  \(e : \operatorname{Spec} k \to X\) a section of the structure
  morphism, and \(x = e(*)\). Then the dual-number points of \(X\) at
  \(x\) over \(\operatorname{Spec} k\) form the \(\kappa(x)\)-linear
  dual of the cotangent space:
  \[
    \{\, g : \operatorname{Spec} k_\varepsilon \to X
        \text{ over } \operatorname{Spec} k
        \mid g(\mathfrak m) = x \,\}
    \;\cong\;
    \operatorname{Hom}_{\kappa(x)}\!\bigl(\mathfrak m_x/\mathfrak m_x^2,\,
      \kappa(x)\bigr).
  \]
\end{corollary}

\begin{proof}\leanok
  \uses{thm:tangent_over_dual_number_cotangent_dual}
  By \cref{lem:tangent_section_rational} the point \(x\) is
  \(k\)-rational, so
  \cref{thm:tangent_over_dual_number_cotangent_dual} applies.
\end{proof}

\begin{definition}[Identity section of a group scheme]\leanok
  \label{def:tangent_identity_section}
  \dcref{custom}
  \lean{AlgebraicGeometry.GroupScheme.identitySection}
  Let \(k\) be a field and \(G\) a monoid (in particular, group)
  scheme over \(k\): a scheme over \(\operatorname{Spec} k\) equipped
  with a monoid-object structure in the category of schemes over
  \(\operatorname{Spec} k\). The \emph{identity section}
  \(e : \operatorname{Spec} k \to G\) is the underlying scheme
  morphism of the unit \(\eta : \mathbf 1 \to G\), where the monoidal
  unit \(\mathbf 1\) is \(\operatorname{Spec} k\) with structure
  morphism the identity.
\end{definition}

\begin{lemma}[The identity section is a section]\leanok
  \label{lem:tangent_identity_section_comp}
  \dcref{custom}
  \lean{AlgebraicGeometry.GroupScheme.identitySection_comp}
  \uses{def:tangent_identity_section}
  In the situation of \cref{def:tangent_identity_section}, the
  identity section satisfies
  \(f \circ e = \operatorname{id}_{\operatorname{Spec} k}\), where
  \(f : G \to \operatorname{Spec} k\) is the structure morphism.
\end{lemma}

\begin{proof}\leanok
  The unit \(\eta\) is a morphism of schemes over
  \(\operatorname{Spec} k\) out of the monoidal unit
  \((\operatorname{Spec} k, \operatorname{id})\); its compatibility
  triangle over \(\operatorname{Spec} k\) is precisely the equation
  \(f \circ e = \operatorname{id}\).
\end{proof}

\begin{theorem}[Tangent space of a group scheme at the identity]\leanok
  \label{thm:tangent_identity_cotangent_dual}
  \dcref{custom}
  \lean{AlgebraicGeometry.GroupScheme.identityDualNumberEquivCotangentSpaceDual}
  \uses{def:tangent_identity_section}
  Let \(k\) be a field and \(G\) a monoid (in particular, group)
  scheme over \(k\), with identity section
  \(e : \operatorname{Spec} k \to G\)
  (\cref{def:tangent_identity_section}) and \(x = e(*)\). Then the
  dual-number points of \(G\) at \(x\) over \(\operatorname{Spec} k\)
  — the Zariski tangent space of \(G\) at the identity in its
  functor-of-points form — form the \(\kappa(x)\)-linear dual of the
  cotangent space \(\mathfrak m_x/\mathfrak m_x^2\).
\end{theorem}

\begin{proof}\leanok
  \uses{cor:tangent_section_cotangent_dual,
        lem:tangent_identity_section_comp}
  By \cref{lem:tangent_identity_section_comp} the identity section is
  a section of the structure morphism, so
  \cref{cor:tangent_section_cotangent_dual} applies at \(x = e(*)\).
\end{proof}

\begin{theorem}[\(\Pic^0_{C/k}\) is a \(k\)-group scheme]\leanok
  \label{thm:pic0_group_structure}
  \dcref{custom}
  \lean{AlgebraicGeometry.Scheme.Pic0.grpObj}
  \uses{def:pic_zero_subscheme,
        thm:identity_component_is_subgroup_homomorphism}
  Let \(C/k\) be a smooth proper geometrically integral curve over a
  field \(k\). The identity component \(\Pic^0_{C/k}\)
  (\cref{def:pic_zero_subscheme}) carries a \(k\)-group-scheme
  structure, inherited from the group scheme \(\Pic_{C/k}\) via the
  clopen inclusion \(\Pic^0_{C/k} \hookrightarrow \Pic_{C/k}\).
\end{theorem}

\begin{proof}\leanok
  \uses{thm:identity_component_is_subgroup_homomorphism,
        thm:pic_is_group_scheme,
        thm:fga_pic_representability}
  The Picard scheme \(\Pic_{C/k}\) is a \(k\)-group scheme
  (\cref{thm:pic_is_group_scheme}) locally of finite type over \(k\)
  (\cref{thm:fga_pic_representability}), and \(\Pic^0_{C/k}\) is by
  definition its identity component. By
  \cref{thm:identity_component_is_subgroup_homomorphism} the identity
  component of a \(k\)-group scheme locally of finite type is a
  subgroup scheme; in particular \(\Pic^0_{C/k}\) inherits a
  \(k\)-group-scheme structure compatible with the inclusion.
\end{proof}

\begin{theorem}[Dual-number points of \(\Pic^0_{C/k}\) at the identity]\leanok
  \label{thm:pic0_tangent_cotangent_dual}
  \dcref{custom}
  \lean{AlgebraicGeometry.Scheme.Pic0.tangentSpaceCotangentDual}
  \uses{def:pic_zero_subscheme,
        thm:pic0_group_structure}
  Let \(C/k\) be a smooth proper geometrically integral curve over a
  field \(k\), and equip \(\Pic^0_{C/k}\)
  (\cref{def:pic_zero_subscheme}) with its \(k\)-group-scheme
  structure (\cref{thm:pic0_group_structure}). Then there
  is an identity section \(e : \operatorname{Spec} k \to
  \Pic^0_{C/k}\) of the structure morphism such that the dual-number
  points of \(\Pic^0_{C/k}\) at \(e(*)\) over \(\operatorname{Spec} k\)
  form the \(\kappa(e(*))\)-linear dual of the cotangent space:
  \[
    \tu T_0 \Pic^0_{C/k}
    \;=\;
    \{\, g : \operatorname{Spec} k_\varepsilon \to \Pic^0_{C/k}
        \text{ over } \operatorname{Spec} k
        \mid g(\mathfrak m) = e(*) \,\}
    \;\cong\;
    \operatorname{Hom}_{\kappa(e(*))}\!\bigl(
      \mathfrak m_{e(*)}/\mathfrak m_{e(*)}^2,\, \kappa(e(*))\bigr).
  \]
  This is the geometric half of \cref{thm:pic0_tangent_space_iso};
  what remains is the identification of the right-hand side with
  \(H^1(C, \mathcal O_C)\).
\end{theorem}

\begin{proof}\leanok
  \uses{thm:tangent_identity_cotangent_dual, thm:pic0_group_structure}
  The group-scheme structure (\cref{thm:pic0_group_structure})
  provides the identity section, and
  \cref{thm:tangent_identity_cotangent_dual} applied to
  \(G = \Pic^0_{C/k}\) gives the identification.
\end{proof}


\begin{theorem}[Tangent space at the identity: \(\tu T_0 \Pic^0_{C/k} \cong H^1(C, \mathcal{O}_C)\)]\leanok
  \label{thm:pic0_tangent_space_iso}
  \dcref{custom}
  \lean{AlgebraicGeometry.Scheme.Pic0.tangentSpaceIso}
  \uses{def:pic_zero_subscheme,
        thm:identity_component_open_subgroup,
        def:genus}

  % SOURCE: \cite{kleiman-picard}, ``The Picard scheme'', \S 5, Thm.~thm:tgtsp;
  % arXiv:math/0504020, p.~41 (read from
  % references/kleiman-picard-src/kleiman-picard.tex, L3265-L3270).
  % SOURCE QUOTE:
  % "\\label{thm:tgtsp}
  %  Assume \(S\) is the spectrum of a field \(k\).  Assume \(\IPic_{X/k}\) exists
  % and represents \(\Pic_{(X/k)\et}\).  Let \(\tu T_0\IPic_{X/k}\) denote the
  % tangent space at \(0\).  Then
  %    \(\)\tu T_0\IPic_{X/k}=\tu H^1(\mc O_{\!X}).\(\)
  %  \"
  \textit{Source: \cite{kleiman-picard}, ``The Picard scheme'', \S 5,
  Thm.~5.11; cf.\ \cite{hartshorne-ag}, ``Algebraic Geometry'', GTM 52,
  IV.1.1 (genus def.\ \(g(C) = \dim_k H^1(C, \mathcal{O}_C)\)).}
  Let \(C/k\) be a smooth proper geometrically integral curve of genus
  \(g = g(C)\) over a field \(k\), and let \(\Pic^0_{C/k}\)
  (\cref{def:pic_zero_subscheme}) be the identity component of the
  Picard scheme of \(C/k\). There is a canonical isomorphism of
  \(k\)-vector spaces
  \[
    \tu T_0 \Pic^0_{C/k} \;\cong\; H^1(C, \mathcal{O}_C).
  \]
  In particular, the \(k\)-dimension of \(\tu T_0 \Pic^0_{C/k}\) equals
  the genus \(g(C) = \dim_k H^1(C, \mathcal{O}_C)\)
  (\cref{def:genus}).
\end{theorem}

% SOURCE QUOTE PROOF: (from references/kleiman-picard-src/kleiman-picard.tex,
% L3358-L3367 + L3404-L3407: the construction of the comparison map and
% its identification as an isomorphism. Kleiman's full proof spans
% L3271-L3408 and is too long for one verbatim block, so we reproduce
% the two paragraphs that contain the load-bearing exact sequence and
% the conclusion; see Notes for Plan Agent.)
% "To compute this kernel, set \(X_\ve:=X\ox_kk_\ve\), and form the truncated
% exponential sequence of sheaves of Abelian groups:
%    \(\)0\to \mc O_{\!X} \to \mc O_{\!X_\ve}^* \to \mc O_{\!X}^* \to 1,\(\)
%  where the first map takes a local section \(b\) to \(1+b\ve\).  This
% sequence is split by the map \(\mc O_{\!X}^*\to \mc O_{\!X_\ve}^*\) defined by
% \(a\mapsto a+0\cdot\ve\).  Hence taking cohomology yields this split exact
% sequence of Abelian groups:
%    \(\)0\to \tu H^1(\mc O_{\!X}) \to \tu H^1(\mc O_{\!X_\ve}^*)
%     \to \tu H^1(\mc O_{\!X}^*) \to 1.\(\)"
% "In Square (\ref{CD:tgtsp}), the two vertical maps are isomorphisms by
% Exercise~\ref{ex:Alr} since \(\IPic_{X/k}\) represents \(\Pic_{(X/k)\et}\) and
% \(k\) is algebraically closed.  Therefore, \(v\)
% too is an isomorphism, as desired."
\begin{proof}
  \uses{thm:fga_pic_representability,
        thm:identity_component_open_subgroup,
        lem:tangent_derivation_cotangent_dual,
        cor:tangent_dual_number_cotangent_dual,
        cor:tangent_spec_dual_number_at,
        thm:pic0_tangent_cotangent_dual,
        thm:dual_number_units_split,
        lem:dual_number_trunc_exp_exact,
        def:dual_number_map}
  We follow Kleiman~\S 5 Thm.~5.11. By
  \cref{thm:fga_pic_representability}, the Picard scheme
  \(\Pic_{C/k}\) exists and represents the étale-sheafified relative
  Picard functor \(\Pic_{(C/k)\et}\); the identity component
  \(\Pic^0_{C/k}\) is an open subscheme of \(\Pic_{C/k}\)
  (\cref{thm:identity_component_open_subgroup}), so the tangent
  spaces at the identity of \(\Pic^0_{C/k}\) and of \(\Pic_{C/k}\)
  agree; we may therefore work with \(\Pic_{C/k}\) throughout.

  \emph{Tangent space via dual numbers.} Let
  \(k_\varepsilon := k[\varepsilon]/(\varepsilon^2)\) be the ring of
  dual numbers. For any \(k\)-scheme \(P\) locally of finite type and
  \(k\)-rational point \(e \in P(k)\), the Zariski tangent space
  \(\tu T_e P\) is in functorial bijection with the set
  \(P(k_\varepsilon)_e\) of \(k_\varepsilon\)-valued points whose
  underlying \(k\)-point is \(e\): such points are, by
  \cref{cor:tangent_spec_dual_number_at}, the local
  \(k\)-algebra maps \(\mathcal O_{P,e} \to k_\varepsilon\), which by
  \cref{cor:tangent_dual_number_cotangent_dual} (via
  \cref{lem:tangent_derivation_cotangent_dual}) form the dual of the
  cotangent space
  \(\mathfrak m_e/\mathfrak m_e^2\). When \(P\) is a \(k\)-group scheme,
  the natural homomorphism \(k_\varepsilon \to k\) (\(\varepsilon
  \mapsto 0\)) induces a group homomorphism
  \(P(k_\varepsilon) \to P(k)\) whose kernel coincides with
  \(P(k_\varepsilon)_e\); under this identification the vector-space
  addition on \(\tu T_e P\) corresponds to the group multiplication on
  the kernel. Specialising to \(P = \Pic_{C/k}\),
  \[
    \tu T_0 \Pic_{C/k}
    \;=\; \ker\bigl(\Pic_{C/k}(k_\varepsilon) \to \Pic_{C/k}(k)\bigr).
  \]
  The identification of \(\tu T_0 \Pic^0_{C/k}\) — the dual-number
  points over the identity section — with the
  \(\kappa(0)\)-linear dual of the cotangent space
  \(\mathfrak m_0/\mathfrak m_0^2\) is
  \cref{thm:pic0_tangent_cotangent_dual}.

  \emph{Computing the kernel via the truncated exponential sequence.}
  Set \(C_\varepsilon := C \otimes_k k_\varepsilon\). On the étale
  site of \(C\) the truncated exponential sequence of sheaves of
  abelian groups
  \[
    0 \to \mathcal{O}_C \to \mathcal{O}_{C_\varepsilon}^{\times}
        \to \mathcal{O}_C^{\times} \to 1,
  \]
  whose first map sends a local section \(b\) to \(1 + b\varepsilon\),
  is split by \(a \mapsto a + 0 \cdot \varepsilon\). Taking sheaf
  cohomology produces a split exact sequence
  \[
    0 \to H^1(C, \mathcal{O}_C)
      \to H^1(C, \mathcal{O}_{C_\varepsilon}^{\times})
      \to H^1(C, \mathcal{O}_C^{\times}) \to 1.
  \]

  \emph{Identification with the tangent space.} Because \(\Pic_{C/k}\)
  represents the étale-sheafification of
  \(T \mapsto H^1(C_T, \mathcal{O}_{C_T}^{\times})\), there is a
  commutative diagram of abelian groups
  \[
    \begin{array}{ccc}
      H^1(C, \mathcal{O}_{C_\varepsilon}^{\times}) &
        \longrightarrow & H^1(C, \mathcal{O}_C^{\times}) \\
      \downarrow & & \downarrow \\
      \Pic_{C/k}(k_\varepsilon) & \longrightarrow &
        \Pic_{C/k}(k),
    \end{array}
  \]
  whose vertical maps are isomorphisms after base change to the
  algebraic closure \(\bar k\) (Kleiman~\S 2, Exercise~2.3;
  invocation legitimised by
  \cref{thm:identity_component_base_change_commutes}). Taking horizontal
  kernels yields a \(k\)-linear comparison map
  \(v : H^1(C, \mathcal{O}_C) \to \tu T_0 \Pic_{C/k}\) which is an
  isomorphism after base change to \(\bar k\), hence already an
  isomorphism over \(k\) (the source and target are finite-dimensional
  \(k\)-vector spaces, and an isomorphism of \(k\)-vector spaces is
  detected after a faithfully flat extension).

  \emph{Compatibility with the additive structure.} The diagram of the
  preceding paragraph is functorial in \(k\)-linear endomorphisms of
  \(k_\varepsilon\); in particular scalar multiplication by \(a \in
  k\) on the dual numbers corresponds, on
  \(H^1(C, \mathcal{O}_C)\), to scalar multiplication by \(a\) in the
  natural \(k\)-module structure. Hence \(v\) is a homomorphism of
  \(k\)-vector spaces, as required.

  Finally, by the genus definition (Hartshorne~IV.1.1,
  \cref{def:genus}), \(\dim_k H^1(C, \mathcal{O}_C) = g(C)\); thus
  \(\dim_k \tu T_0 \Pic^0_{C/k} = g(C)\).
\end{proof}

\subsection{The truncated-exponential splitting, ring level}
\label{sec:pic0_trunc_exp}

The truncated-exponential sequence of sheaves used in the proof of
\cref{thm:pic0_tangent_space_iso} is, on each ring of sections, the
following split short exact sequence of groups attached to the dual
numbers \(R[\varepsilon] := R[\varepsilon]/(\varepsilon^2)\) of a
commutative ring \(R\):
\[
  1 \longrightarrow (R, +)
    \xrightarrow{\ b \mapsto 1 + b\varepsilon\ }
    (R[\varepsilon])^{\times}
    \xrightarrow{\ \varepsilon \mapsto 0\ }
    R^{\times} \longrightarrow 1,
\]
split by \(a \mapsto a + 0\cdot\varepsilon\). This subsection develops
the sequence and its naturality in \(R\); the sheaf-level upgrade on the
dual-number thickening \(C_\varepsilon\) is the remaining input to
\cref{thm:pic0_tangent_space_iso}.

\begin{lemma}[Reduction mod \(\varepsilon\) is a ring homomorphism]\leanok
  \label{lem:dual_number_fst_ring_hom}
  \dcref{custom}
  \lean{DualNumber.fstRingHom}
  Let \(R\) be a commutative ring. The first-component projection
  \(a + b\varepsilon \mapsto a\) is a ring homomorphism
  \(R[\varepsilon] \to R\).
\end{lemma}

\begin{proof}\leanok
  The projection is additive and unital, and multiplicative because
  \((a + b\varepsilon)(a' + b'\varepsilon)
  = aa' + (ab' + a'b)\varepsilon\) has first component \(aa'\).
\end{proof}

\begin{lemma}[The truncated exponential is a unit]\leanok
  \label{lem:dual_number_trunc_exp_unit}
  \dcref{custom}
  \lean{DualNumber.truncExpUnit}
  Let \(R\) be a commutative ring and \(b \in R\). Then
  \(1 + b\varepsilon\) is a unit of \(R[\varepsilon]\), with inverse
  \(1 - b\varepsilon\).
\end{lemma}

\begin{proof}\leanok
  \((1 + b\varepsilon)(1 - b\varepsilon) = 1 - b^2\varepsilon^2 = 1\)
  since \(\varepsilon^2 = 0\).
\end{proof}

\begin{definition}[The truncated exponential]\leanok
  \label{def:dual_number_trunc_exp}
  \dcref{custom}
  \lean{DualNumber.truncExp}
  \uses{lem:dual_number_trunc_exp_unit}
  Let \(R\) be a commutative ring. The \emph{truncated exponential} is
  the group homomorphism
  \((R, +) \to (R[\varepsilon])^{\times}\), \(b \mapsto 1 + b\varepsilon\);
  it is a homomorphism because
  \((1 + b\varepsilon)(1 + c\varepsilon) = 1 + (b + c)\varepsilon\)
  (\(\varepsilon^2 = 0\)).
\end{definition}

\begin{lemma}[The truncated exponential is injective]\leanok
  \label{lem:dual_number_trunc_exp_injective}
  \dcref{custom}
  \lean{DualNumber.truncExp_injective}
  \uses{def:dual_number_trunc_exp}
  The truncated exponential \(b \mapsto 1 + b\varepsilon\) is injective.
\end{lemma}

\begin{proof}\leanok
  The \(\varepsilon\)-component of \(1 + b\varepsilon\) is \(b\).
\end{proof}

\begin{definition}[Reduction mod \(\varepsilon\) on unit groups]\leanok
  \label{def:dual_number_units_fst}
  \dcref{custom}
  \lean{DualNumber.unitsFst}
  \uses{lem:dual_number_fst_ring_hom}
  Let \(R\) be a commutative ring. Reduction mod \(\varepsilon\)
  (\cref{lem:dual_number_fst_ring_hom}) induces a group homomorphism
  \((R[\varepsilon])^{\times} \to R^{\times}\) on unit groups.
\end{definition}

\begin{definition}[The constant section of the unit groups]\leanok
  \label{def:dual_number_units_inl}
  \dcref{custom}
  \lean{DualNumber.unitsInl}
  Let \(R\) be a commutative ring. The inclusion
  \(a \mapsto a + 0 \cdot \varepsilon\) (a ring homomorphism
  \(R \to R[\varepsilon]\)) induces a group homomorphism
  \(R^{\times} \to (R[\varepsilon])^{\times}\) on unit groups.
\end{definition}

\begin{lemma}[The unit-group sequence is split]\leanok
  \label{lem:dual_number_units_fst_split}
  \dcref{custom}
  \lean{DualNumber.unitsFst_comp_unitsInl}
  \uses{def:dual_number_units_fst, def:dual_number_units_inl}
  The composite
  \(R^{\times} \to (R[\varepsilon])^{\times} \to R^{\times}\)
  of \cref{def:dual_number_units_inl} followed by
  \cref{def:dual_number_units_fst} is the identity. In particular
  reduction mod \(\varepsilon\) on unit groups is surjective.
\end{lemma}

\begin{proof}\leanok
  The first component of \(a + 0\cdot\varepsilon\) is \(a\).
\end{proof}

\begin{lemma}[First component of an inverse unit]\leanok
  \label{lem:dual_number_fst_unit_inv}
  \dcref{custom}
  \lean{DualNumber.fst_inv_mul_fst}
  \uses{lem:dual_number_fst_ring_hom}
  For a unit \(u \in (R[\varepsilon])^{\times}\), the first components
  of \(u^{-1}\) and of \(u\) multiply to \(1\); that is, reduction mod
  \(\varepsilon\) carries \(u^{-1}\) to the inverse of the reduction of
  \(u\).
\end{lemma}

\begin{proof}\leanok
  Apply the ring homomorphism of
  \cref{lem:dual_number_fst_ring_hom} to \(u^{-1} u = 1\).
\end{proof}

\begin{lemma}[Exactness of the truncated-exponential sequence]\leanok
  \label{lem:dual_number_trunc_exp_exact}
  \dcref{custom}
  \lean{DualNumber.truncExp_range_eq_ker_unitsFst}
  \uses{def:dual_number_trunc_exp, def:dual_number_units_fst}
  The range of the truncated exponential
  \((R,+) \to (R[\varepsilon])^{\times}\) equals the kernel of
  reduction mod \(\varepsilon\) on unit groups: a unit of
  \(R[\varepsilon]\) reduces to \(1\) if and only if it is of the form
  \(1 + b\varepsilon\) for a (unique) \(b \in R\).
\end{lemma}

\begin{proof}\leanok
  \uses{lem:dual_number_trunc_exp_injective}
  A unit \(1 + b\varepsilon\) visibly reduces to \(1\). Conversely a
  unit \(u = a + m\varepsilon\) with \(a = 1\) equals
  \(1 + m\varepsilon\), the truncated exponential of \(m\);
  uniqueness of \(m\) is
  \cref{lem:dual_number_trunc_exp_injective}.
\end{proof}

\begin{theorem}[Splitting of the dual-number unit group]\leanok
  \label{thm:dual_number_units_split}
  \dcref{custom}
  \lean{DualNumber.unitsEquivProd}
  \uses{def:dual_number_trunc_exp, def:dual_number_units_fst,
        def:dual_number_units_inl}
  Let \(R\) be a commutative ring. There is a group isomorphism
  \[
    (R[\varepsilon])^{\times} \;\cong\; R^{\times} \times (R, +),
  \]
  sending a unit \(u = a + m\varepsilon\) (so \(a \in R^{\times}\)) to
  \((a, a^{-1}m)\), with inverse
  \((a, b) \mapsto a(1 + b\varepsilon) = a + ab\,\varepsilon\).
\end{theorem}

\begin{proof}\leanok
  \uses{lem:dual_number_fst_unit_inv, lem:dual_number_units_fst_split,
        lem:dual_number_trunc_exp_exact}
  The two maps are mutually inverse: composing them in either order and
  comparing components uses only
  \(a^{-1}a = 1\) (\cref{lem:dual_number_fst_unit_inv}) and
  \(\varepsilon^2 = 0\). Multiplicativity of the forward map: for
  units \(u = a + m\varepsilon\) and \(v = a' + m'\varepsilon\),
  \(uv = aa' + (am' + a'm)\varepsilon\), and
  \((aa')^{-1}(am' + a'm) = a'^{-1}m' + a^{-1}m\), which in the target
  \(R^{\times} \times (R,+)\) is exactly the product of the images
  \((a, a^{-1}m)\) and \((a', a'^{-1}m')\); commutativity of \(R\) is
  used here.
\end{proof}

\begin{definition}[Functoriality of the dual numbers]\leanok
  \label{def:dual_number_map}
  \dcref{custom}
  \lean{DualNumber.mapRingHom}
  A ring homomorphism \(f : R \to S\) of commutative rings induces a
  ring homomorphism
  \(R[\varepsilon] \to S[\varepsilon]\),
  \(a + b\varepsilon \mapsto f(a) + f(b)\varepsilon\).
\end{definition}

\begin{lemma}[Functoriality: identity]\leanok
  \label{lem:dual_number_map_id}
  \dcref{custom}
  \lean{DualNumber.mapRingHom_id}
  \uses{def:dual_number_map}
  The ring homomorphism \(R[\varepsilon] \to R[\varepsilon]\) induced
  by the identity of \(R\) is the identity.
\end{lemma}

\begin{proof}\leanok
  Immediate from the componentwise formula.
\end{proof}

\begin{lemma}[Functoriality: composition]\leanok
  \label{lem:dual_number_map_comp}
  \dcref{custom}
  \lean{DualNumber.mapRingHom_comp}
  \uses{def:dual_number_map}
  For ring homomorphisms \(f : R \to S\) and \(g : S \to T\) of
  commutative rings, the map on dual numbers induced by
  \(g \circ f\) is the composite of the maps induced by \(f\) and by
  \(g\).
\end{lemma}

\begin{proof}\leanok
  Immediate from the componentwise formula.
\end{proof}

\begin{lemma}[Naturality of the truncated exponential]\leanok
  \label{lem:dual_number_map_trunc_exp}
  \dcref{custom}
  \lean{DualNumber.map_mapRingHom_truncExpUnit}
  \uses{def:dual_number_map, lem:dual_number_trunc_exp_unit}
  For a ring homomorphism \(f : R \to S\) of commutative rings, the
  induced map on dual-number units carries \(1 + b\varepsilon\) to
  \(1 + f(b)\varepsilon\).
\end{lemma}

\begin{proof}\leanok
  Immediate from the componentwise formula.
\end{proof}

\begin{lemma}[Naturality of reduction mod \(\varepsilon\)]\leanok
  \label{lem:dual_number_units_fst_natural}
  \dcref{custom}
  \lean{DualNumber.unitsFst_map_mapRingHom}
  \uses{def:dual_number_map, def:dual_number_units_fst}
  For a ring homomorphism \(f : R \to S\) of commutative rings, the
  square
  \[
    \begin{array}{ccc}
      (R[\varepsilon])^{\times} & \longrightarrow &
        (S[\varepsilon])^{\times} \\
      \downarrow & & \downarrow \\
      R^{\times} & \longrightarrow & S^{\times}
    \end{array}
  \]
  of reduction mod \(\varepsilon\) on unit groups against the maps
  induced by \(f\) commutes.
\end{lemma}

\begin{proof}\leanok
  Both composites send \(a + b\varepsilon\) to \(f(a)\), on first
  components.
\end{proof}

\subsection{Chart triviality over a nilpotent thickening}

An invertible sheaf on \(C \times_k \Spec k[\varepsilon]\) whose restriction
along \(\varepsilon \mapsto 0\) is trivial is trivial on each base-changed
affine chart. Over an affine chart \(V = \Spec A\) the base change is
\(\Spec A[\varepsilon]\) and an invertible sheaf is an invertible module, so
the statement is the following piece of commutative algebra, in which nothing
is special to the dual numbers beyond nilpotency of the augmentation ideal.

\begin{lemma}[An invertible module generated by one element is free]\leanok
  \label{lem:invertible_cyclic_free}
  \dcref{custom}
  \lean{Module.Invertible.free_of_span_singleton_eq_top}
  Let \(R\) be a commutative ring and \(M\) an invertible \(R\)-module. If
  \(M\) is generated by a single element \(m\), then \(M \cong R\).
\end{lemma}

\begin{proof}\leanok
  The map \(R \to M\), \(r \mapsto rm\), is surjective precisely because \(m\)
  generates. For an invertible module a surjection onto it from an invertible
  module is automatically injective, and \(R\) is invertible over itself;
  hence the map is an isomorphism.

  The hypothesis of invertibility cannot be dropped: \(\mathbb{Z}/2\) is
  generated by one element over \(\mathbb{Z}\) and is not free.
\end{proof}

\begin{lemma}[Nakayama for a nilpotent ideal]\leanok
  \label{lem:nilpotent_nakayama}
  \dcref{custom}
  \lean{Submodule.top_le_span_sup_pow_smul_top}
  Let \(I\) be an ideal of a commutative ring \(R\), let \(M\) be an
  \(R\)-module and let \(N \subseteq M\) be a submodule with
  \(N + IM = M\). Then \(N + I^{n}M = M\) for every \(n \ge 0\).
\end{lemma}

\begin{proof}\leanok
  Induction on \(n\). Substituting \(M = N + IM\) into \(I^{n}M\) gives
  \(I^{n}M = I^{n}N + I^{n+1}M\); the first summand lies in \(N\) and the
  second is \(I^{n+1}M\), so \(N + I^{n}M \subseteq N + I^{n+1}M\).

  No finiteness hypothesis on \(N\) or \(M\) is required, and no hypothesis
  relating \(I\) to the Jacobson radical. This matters for the application:
  the modules of interest are section rings of affine charts of a curve, which
  are not finite over the base field.
\end{proof}

\begin{proposition}[Triviality over a nilpotent thickening]\leanok
  \label{prp:invertible_free_of_nilpotent}
  \dcref{custom}
  \lean{Module.Invertible.free_of_nilpotent_of_exists_sub_smul_mem}
  \uses{lem:invertible_cyclic_free, lem:nilpotent_nakayama}
  Let \(I\) be a nilpotent ideal of a commutative ring \(R\) and let \(M\) be
  an invertible \(R\)-module. If the reduction \(M/IM\) is generated by the
  image of a single element \(m \in M\), then \(M \cong R\).
\end{proposition}

\begin{proof}\leanok
  The hypothesis says \(Rm + IM = M\). By
  \cref{lem:nilpotent_nakayama}, \(Rm + I^{n}M = M\) for all \(n\), and
  \(I^{j} = 0\) for some \(j\), so \(Rm = M\). Now apply
  \cref{lem:invertible_cyclic_free}.
\end{proof}

\begin{lemma}[The augmentation ideal is square-zero]\leanok
  \label{lem:dual_number_aug_square_zero}
  \dcref{custom}
  \lean{DualNumber.augIdeal_mul_self_eq_bot}
  The kernel of the augmentation \(A[\varepsilon] \to A\),
  \(a + b\varepsilon \mapsto a\), satisfies \(\mathfrak{a}^{2} = 0\), and
  \(\mathfrak{a} = (\varepsilon)\).
\end{lemma}

\begin{proof}\leanok
  The kernel consists of the elements with vanishing constant term. For two
  such elements the product has vanishing constant term as well, and its
  \(\varepsilon\)-coefficient is
  \(a_{0}b_{1} + b_{0}a_{1} = 0\) since \(a_{0} = b_{0} = 0\). An element
  \(b\varepsilon\) of the kernel equals \(\varepsilon \cdot b\), whence the
  second assertion.
\end{proof}

\begin{corollary}[Chart triviality]\leanok
  \label{cor:dual_number_chart_triviality}
  \dcref{custom}
  \lean{DualNumber.free_of_cyclic_mod_eps}
  \uses{prp:invertible_free_of_nilpotent, lem:dual_number_aug_square_zero}
  Let \(A\) be a commutative ring and \(M\) an invertible
  \(A[\varepsilon]\)-module whose restriction along \(\varepsilon \mapsto 0\)
  is generated by the image of a single element. Then \(M \cong
  A[\varepsilon]\).
\end{corollary}

\begin{proof}\leanok
  The augmentation ideal is square-zero by
  \cref{lem:dual_number_aug_square_zero}, hence nilpotent; apply
  \cref{prp:invertible_free_of_nilpotent}.
\end{proof}

\section{Smoothness of \(\Pic^0_{C/k}\)}
\label{sec:pic0_smooth}

By the tangent-space dimension bound for a Noetherian local ring,
\(\dim_0 \Pic^0_{C/k} \le \dim_k \tu T_0 \Pic^0_{C/k} = g\), with
equality if and only if \(\Pic^0_{C/k}\) is regular at \(0\). For a
smooth proper curve, this inequality becomes an equality: smoothness
at the identity propagates to smoothness everywhere by
translation, hence \(\Pic^0_{C/k}\) is a smooth \(k\)-group scheme of
dimension \(g\).

\begin{theorem}[Smoothness of \(\Pic^0_{C/k}\)]
  \label{thm:pic0_smooth}
  \dcref{custom}
  \lean{AlgebraicGeometry.Scheme.Pic0.smooth}
  \uses{thm:pic0_tangent_space_iso,
        def:pic_zero_subscheme,
        thm:identity_component_is_subgroup_homomorphism}

  % SOURCE: \cite{kleiman-picard}, ``The Picard scheme'', \S 5, Cor.~cor:sm + Cor.~cor:ch0;
  % arXiv:math/0504020, p.~44 (read from
  % references/kleiman-picard-src/kleiman-picard.tex, L3421-L3427 +
  % L3442-L3446).
  % SOURCE QUOTE:
  % "\\label{cor:sm} Assume \(S\) is the spectrum of a field
  % \(k\). Assume \(\IPic_{X/k}\) exists and represents \(\Pic_{(X/k)\et}\).  Then
  %    \(\)\dim \IPic_{X/k}\le \dim\tu H^1(\mc O_{\!X}).\(\)
  %  Equality holds if and only if \(\IPic_{X/k}\) is smooth at \(0;\) if so, then
  % \(\IPic_{X/k}\) is smooth of dimension \(\dim\tu H^1(\mc O_{\!X})\)
  % everywhere.
  %  \"
  % SOURCE QUOTE:
  % "\\label{cor:ch0}
  %  Assume \(S\) is the spectrum of a field \(k\). Assume \(\IPic_{X/k}\) exists
  % and represents \(\Pic_{(X/k)\et}\). If \(k\) is of characteristic \(0\), then
  % \(\IPic_{X/k}\) is smooth of dimension \(\dim\tu H^1(\mc O_{\!X})\) everywhere.
  %  \"
  \textit{Source: \cite{kleiman-picard}, ``The Picard scheme'', \S 5,
  Cor.~5.13 + Cor.~5.14; cf.\ \cite{milne-av}, ``Abelian Varieties''
  (course notes, 2008), \S III.1, Rmk.~1.4(e), p.~86 (dimension of \(J\)
  equals the genus).}
  Let \(C/k\) be a smooth proper geometrically integral curve of genus
  \(g\) over a field \(k\). The identity component \(\Pic^0_{C/k}\)
  (\cref{def:pic_zero_subscheme}) is smooth over \(k\) of dimension
  \(g\): the structural morphism
  \((\Pic^0_{C/k}) \to \Spec k\) is smooth, and every closed point of
  \(\Pic^0_{C/k}\) lies in a smooth open neighbourhood.
\end{theorem}

% SOURCE QUOTE PROOF: (from references/kleiman-picard-src/kleiman-picard.tex,
% L3428-L3439: the proof of Cor.~cor:sm. We reproduce the full proof,
% which deduces smoothness everywhere from smoothness at \(0\) by
% translation.)
% "\begin{proof} Plainly we may assume \(k\) is algebraically closed. Then,
% given any closed point \(\lambda\in\IPic_{X/k}\), there is an automorphism
% of \(\IPic_{X/k}\) that carries 0 to \(\lambda\), namely, ``multiplication''
% by \(\lambda\).  So \(\IPic_{X/k}\) has the same dimension at \(\lambda\) as
% at 0, and \(\IPic_{X/k}\) is smooth at \(\lambda\) if and only if it is
% smooth at 0.
%
% By general principles, \(\dim_0 \IPic_{X/k}\le \dim\tu T_0\IPic_{X/k}\),
% and equality holds if and only if \(\IPic_{X/k}\) is regular at \(0\).
% Moreover, \(\IPic_{X/k}\) is regular at \(0\) if and only if it is smooth
% at \(0\) since \(k\) is algebraically closed.  Therefore, the corollary
% results from Theorem~\ref{thm:tgtsp}.
%  \end{proof}"
\begin{proof}
  \uses{thm:pic0_tangent_space_iso,
        thm:identity_component_is_subgroup_homomorphism,
        thm:identity_component_base_change_commutes}
  We follow Kleiman~\S 5 Cor.~5.13, with the curve-specific dimension
  input supplied by \cref{thm:pic0_tangent_space_iso}.

  Since smoothness of a scheme over \(k\) descends along faithfully
  flat extensions, we may replace \(k\) by its algebraic closure
  \(\bar k\) (using \cref{thm:identity_component_base_change_commutes}
  to identify \((\Pic^0_{C/k})_{\bar k}\) with \(\Pic^0_{C_{\bar k}/\bar k}\)).
  Over \(\bar k\) every closed point of \(\Pic^0_{C/k}\) is
  \(\bar k\)-rational, and the group law translates the identity \(0\)
  onto an arbitrary closed point: the translation morphism
  \(t_\lambda : \Pic^0_{C/k} \to \Pic^0_{C/k}\), \(x \mapsto x +
  \lambda\), is an automorphism by
  \cref{thm:identity_component_is_subgroup_homomorphism}. Consequently
  \(\Pic^0_{C/k}\) has the same local dimension at \(\lambda\) as at
  \(0\), and is smooth at \(\lambda\) if and only if it is smooth at
  \(0\).

  By the standard tangent-space dimension bound, for any Noetherian
  local ring \((A, \mathfrak{m})\) with residue field \(k\) one has
  \(\dim A \le \dim_k (\mathfrak{m}/\mathfrak{m}^2)\), with equality
  if and only if \(A\) is regular. Over an algebraically closed field,
  regularity at a \(\bar k\)-rational point is equivalent to
  smoothness. Hence
  \(
    \dim_0 \Pic^0_{C/k} \le \dim_k \tu T_0 \Pic^0_{C/k} = g(C)
  \)
  by \cref{thm:pic0_tangent_space_iso}, with equality iff
  \(\Pic^0_{C/k}\) is regular (equivalently, smooth) at \(0\).

  It remains to show that for \(C/k\) a smooth proper curve, equality
  holds (so smoothness of \(\Pic^0_{C/k}\) follows). For a smooth
  proper curve there is no obstruction: as Kleiman~\S 5 Ex.~5.23
  notes, projectivity and flatness of \(C/k\) together with the
  smoothness of \(C\) over \(k\) imply that \(\Pic^0_{C/k}\) is a
  smooth quasi-projective family of relative dimension \(p_a(C) =
  g(C)\). In characteristic zero this reduces to Cartier's theorem
  (Kleiman~\S 5 Cor.~5.14): every \(k\)-group scheme locally of
  finite type over a field of characteristic zero is smooth; in
  positive characteristic the smoothness uses the curve-specific
  Hodge-theoretic vanishing
  \(H^2(C, \mathcal{O}_C) = 0\) (the source-target of the obstruction
  to lifting tangent vectors to first-order deformations), which holds
  for any curve since \(C\) has dimension \(1\). In either case
  \(\Pic^0_{C/k}\) is smooth at \(0\); by the translation argument
  above, it is smooth everywhere.

  Thus \(\Pic^0_{C/k} \to \Spec k\) is smooth and
  \(\dim \Pic^0_{C/k} = g(C)\), completing the proof.
\end{proof}

\section{Properness of \(\Pic^0_{C/k}\)}
\label{sec:pic0_proper}

A smooth proper geometrically integral curve is geometrically normal,
so Kleiman's \(\Pic^0\)-projectivity theorem applies and upgrades the
quasi-projective conclusion to projective: \(\Pic^0_{C/k}\) is
projective over \(k\), in particular proper.

\begin{theorem}[Properness of \(\Pic^0_{C/k}\)]
  \label{thm:pic0_proper}
  \dcref{custom}
  \lean{AlgebraicGeometry.Scheme.Pic0.proper}
  \uses{def:pic_zero_subscheme,
        thm:identity_component_finite_type_geom_irreducible}

  % SOURCE: \cite{kleiman-picard}, ``The Picard scheme'', \S 5, Thm.~th:qpp\&p;
  % arXiv:math/0504020, p.~36 (read from
  % references/kleiman-picard-src/kleiman-picard.tex, L2935-L2940).
  % SOURCE QUOTE:
  % "\\label{th:qpp&p}
  %  Assume \(S\) is the spectrum of a field \(k\).  Assume \(X/k\) is projective
  % and geometrically integral.  Then \(\IPicz_{X/k}\) exists and is
  % quasi-projective.  If also \(X/k\) is geometrically normal, then
  % \(\IPicz_{X/k}\) is projective.
  %  \"
  \textit{Source: \cite{kleiman-picard}, ``The Picard scheme'', \S 5,
  Thm.~5.4; cf.\ \cite{milne-av}, ``Abelian Varieties'' (course notes,
  2008), \S I.6 (abelian varieties are projective, p.~27).}
  Let \(C/k\) be a smooth proper geometrically integral curve over a
  field \(k\). The identity component \(\Pic^0_{C/k}\)
  (\cref{def:pic_zero_subscheme}) is projective over \(k\), in
  particular proper: the structural morphism
  \((\Pic^0_{C/k}) \to \Spec k\) is a projective morphism, hence
  satisfies the universal closedness and separatedness conditions of
  properness.
\end{theorem}

% SOURCE QUOTE PROOF: (from references/kleiman-picard-src/kleiman-picard.tex,
% L2941-L2984: the proof of th:qpp&p uses the Chevalley--Rosenlicht
% structure theorem to reduce projectivity to ruling out non-constant
% maps from the multiplicative group; the curve case relies on
% geometric normality of \(C\).)
% "Suppose \(X\) is also geometrically normal.  Since \(\IPicz_{X/k}\) is
% quasi-projective, to prove it is projective, it suffices to prove it is
% proper.  By Lemma~\ref{lem:agps}, forming \(\IPicz_{X/k}\) commutes
% with extending \(k\).  And by \cite[2.7.1(vii)]{ega-iv2}, a \(k\)-scheme is
% complete if (and only if) it is after extending \(k\).  So we may, and do,
% assume \(k\) is algebraically closed.
%
% Recall the structure theorem of Chevalley and Rosenlicht for algebraic
% groups, or reduced connected group schemes of finite type over \(k\); see
% \cite[Thm.~1.1, p.~3]{conrad-2002}.  The theorem says that every algebraic
% group is an extension of an Abelian variety (or complete algebraic
% group) by a linear (or affine) algebraic group.  Recall also that every
% solvable linear algebraic \(k\)-group is triangularizable (the
% Lie--Kolchin theorem); so, if it's nontrivial, then it contains a copy
% of the multiplicative group or of the additive group; see \cite[(10.5)
% and (10.2)]{borel-1969}.  Now,  \(\IPic_{X/k}\) is commutative, so solvable.
% Hence it suffices to show that, if \(T\) denotes the
% affine line minus the origin, then every \(k\)-map \(t\:T\to
% (\IPicz_{X/k})_{\red}\) is constant."
\begin{proof}
  \uses{thm:fga_pic_representability,
        thm:identity_component_finite_type_geom_irreducible,
        thm:identity_component_base_change_commutes}
  We follow Kleiman~\S 5 Thm.~5.4, specialised to \(X = C\) a
  smooth proper geometrically integral curve.

  \emph{Quasi-projectivity.} By
  \cref{thm:fga_pic_representability}, \(\Pic_{C/k}\) exists,
  represents \(\Pic_{(C/k)\et}\), and is locally of finite type; its
  identity component \(\Pic^0_{C/k}\) is of finite type by
  \cref{thm:identity_component_finite_type_geom_irreducible}.
  Kleiman's quasi-projectivity statement
  (\S 5 Thm.~5.4, first conclusion) then gives that
  \(\Pic^0_{C/k}\) is a quasi-projective \(k\)-scheme, since \(C/k\)
  is projective and geometrically integral.

  \emph{Reduction to algebraically closed \(k\).} Properness descends
  through faithfully flat extensions of the base field
  (EGA IV\(_2\) 2.7.1(vii)), and formation of \(\Pic^0_{C/k}\)
  commutes with extension of \(k\)
  (\cref{thm:identity_component_base_change_commutes}). Hence it
  suffices to prove that \((\Pic^0_{C/k})_{\bar k} =
  \Pic^0_{C_{\bar k}/\bar k}\) is proper, and we may assume
  \(k = \bar k\) is algebraically closed throughout.

  \emph{Reduction via Chevalley--Rosenlicht.} A quasi-projective
  scheme is projective if and only if it is proper, so it suffices to
  show that \(\Pic^0_{C/k}\) is proper over \(\bar k\). By the
  Chevalley--Rosenlicht structure theorem (Conrad, ``A modern proof of
  Chevalley's theorem on algebraic groups'', Thm.~1.1, p.~3), every
  reduced connected algebraic \(\bar k\)-group is an extension of an
  abelian variety by a linear (affine) algebraic group. The
  multiplicative group \(\mathbb{G}_m\) and the additive group
  \(\mathbb{G}_a\) generate every connected solvable linear group
  (Lie--Kolchin); since \(\Pic^0_{C/k}\) is commutative, hence
  solvable, properness will follow once we show that every
  \(\bar k\)-morphism \(t : \mathbb{G}_m \to (\Pic^0_{C/k})_{\red}\)
  is constant.

  \emph{No non-constant map from \(\mathbb{G}_m\) (uses geometric
  normality of \(C\)).} The smooth proper curve \(C/\bar k\) is
  in particular geometrically normal. Suppose
  \(t : T \to \Pic^0_{C/\bar k}\), where \(T = \mathbb{A}^1
  \setminus \{0\}\) is the multiplicative group; by the comparison
  theorem (Kleiman~\S 2 Thm.~2.5) \(t\) is represented by an
  invertible sheaf \(\mathcal{L}\) on \(C \times T\). Since
  \(C \times T\) is normal and integral, every invertible sheaf is the
  sheaf of a Cartier divisor: write
  \(\mathcal{L} = \mathcal{O}_{C \times T}(D)\) for some divisor \(D\)
  (Hartshorne II Ex.~6.15). Restricting \(\mathcal{L}\) to the generic
  fibre of the projection \(p : C \times T \to C\): the generic fibre
  is \(T\) localised at the generic point of \(C\), and an invertible
  sheaf on an open subscheme of \(\mathbb{A}^1\) is trivial. Hence
  there is a rational function \(\phi\) on \(C \times T\) such that
  \((\phi) + D\) restricts to the trivial divisor on the generic
  fibre. Pick any section \(s : C \to C \times T\) of \(p\) (e.g.\ at
  a closed point of \(T\)), and set \(E := s^*((\phi) + D)\); then
  \(E\) is a divisor on \(C\) with the property that \(p^* E\) and
  \((\phi) + D\) coincide as cycles, hence as divisors (Altman--Kleiman
  AK70, Prp.~3.10, p.~139, valid since \(C \times T\) is normal).
  Therefore \(\mathcal{L} \cong p^* \mathcal{O}_C(E)\), so the
  associated map \(t : T \to \Pic^0_{C/\bar k}\) factors through the
  generic fibre and is constant. This rules out non-constant maps
  \(\mathbb{G}_m \to \Pic^0_{C/k}\); the additive case
  \(\mathbb{G}_a \to \Pic^0_{C/k}\) is analogous (and follows from the
  multiplicative case via the standard embedding
  \(\mathbb{G}_a \hookrightarrow \mathbb{G}_m\) on solvable
  subgroups). Hence the linear part in the Chevalley--Rosenlicht
  decomposition of \(\Pic^0_{C/k}\) is trivial, so \(\Pic^0_{C/k}\) is
  an abelian variety; in particular, it is proper. Combined with
  quasi-projectivity, this yields projectivity of \(\Pic^0_{C/k}\)
  over \(\bar k\), hence over \(k\) by faithfully flat descent of
  properness.
\end{proof}

\section{Geometric irreducibility of \(\Pic^0_{C/k}\)}
\label{sec:pic0_geom_irred}

Geometric irreducibility is delivered directly by the abstract
identity-component substrate
(\cref{thm:identity_component_finite_type_geom_irreducible}) applied
to the \(k\)-group scheme \(G = \Pic_{C/k}\). We restate the
specialisation here for the assembly downstream.

\begin{theorem}[Geometric irreducibility of \(\Pic^0_{C/k}\)]\leanok
  \label{thm:pic0_geom_irred}
  \dcref{custom}
  \lean{AlgebraicGeometry.Scheme.Pic0.geometricallyIrreducible}
  \uses{def:pic_zero_subscheme,
        thm:identity_component_finite_type_geom_irreducible}

  % SOURCE: \cite{kleiman-picard}, ``The Picard scheme'', \S 5, Prp.~prp:pic0;
  % arXiv:math/0504020, p.~36 (read from
  % references/kleiman-picard-src/kleiman-picard.tex, L2921-L2929).
  % SOURCE QUOTE:
  % "\\label{prp:pic0}
  %   Assume \(S\) is the spectrum of a field \(k\).  Assume \(\IPic_{X/k}\)
  % exists and represents \(\Pic_{(X/k)\fppf}\).  Then \(\IPic_{X/k}\) is
  % separated, and it is smooth if it has a geometrically reduced open
  % subscheme.  Furthermore, the connected component of the identity
  % \(\IPicz_{X/k}\) is an open and closed group subscheme of finite type; it
  % is geometrically irreducible; and forming it commutes with extending
  % \(k\).
  %  \"
  \textit{Source: \cite{kleiman-picard}, ``The Picard scheme'', \S 5, Prp.~5.3
  (geometric irreducibility of the connected component of the
  identity).}
  Let \(C/k\) be a smooth proper geometrically integral curve over a
  field \(k\). The identity component \(\Pic^0_{C/k}\)
  (\cref{def:pic_zero_subscheme}) is geometrically irreducible: after
  base change to the algebraic closure \(\bar k\), the scheme
  \((\Pic^0_{C/k})_{\bar k}\) is irreducible.
\end{theorem}

\begin{proof}\leanok
  \uses{thm:identity_component_finite_type_geom_irreducible,
        thm:identity_component_base_change_commutes,
        thm:fga_pic_representability}
  Apply the abstract identity-component substrate to the \(k\)-group
  scheme \(G = \Pic_{C/k}\), which is locally of finite type by
  \cref{thm:fga_pic_representability}. By
  \cref{thm:identity_component_finite_type_geom_irreducible}, the
  identity component \((\Pic_{C/k})^0 = \Pic^0_{C/k}\) is of finite
  type and geometrically irreducible over \(k\). The base-change
  compatibility of
  \cref{thm:identity_component_base_change_commutes} identifies
  \((\Pic^0_{C/k})_{\bar k}\) with \(\Pic^0_{C_{\bar k}/\bar k}\),
  closing the proof.
\end{proof}

\section{Assembly: \(\Pic^0_{C/k}\) is an abelian variety}
\label{sec:pic0_isAbelianVariety}

A \(k\)-group scheme is an \emph{abelian variety} iff its underlying
\(k\)-scheme is a complete (proper) connected group variety
(Milne~\S I.1, p.~8). Combining the smoothness, properness, and
geometric irreducibility statements of this chapter with the
identity-component group-scheme structure of
\cref{chap:Picard_IdentityComponent} packages \(\Pic^0_{C/k}\) as an
abelian variety of dimension \(g\). This is the load-bearing A.3.vii
gate for Route~A.

\begin{theorem}[\(\Pic^0_{C/k}\) is an abelian variety]\leanok
  \label{thm:pic0_isAbelianVariety}
  \dcref{custom}
  \lean{AlgebraicGeometry.Scheme.Pic0.isAbelianVariety}
  \uses{thm:pic0_smooth,
        thm:pic0_proper,
        thm:pic0_geom_irred,
        thm:identity_component_is_subgroup_homomorphism}

  % SOURCE: \cite{milne-av}, ``Abelian Varieties'' (course notes, 2008),
  % \S I.1, p.~8 (definition of abelian variety) (read from
  % references/abelian-varieties.pdf, PDF p.~14).
  % SOURCE QUOTE:
  % "A complete connected group variety is called an abelian variety. As we
  % shall see, they are projective, and (fortunately) commutative. Their
  % group laws will be written additively. Thus \(t_a\) is now denoted
  % \(x \mapsto x + a\) and \(e\) is usually denoted \(0\)."
  % SOURCE: \cite{kleiman-picard}, ``The Picard scheme'', \S 5, Rmk.~rmk:Jac;
  % arXiv:math/0504020, p.~51 (read from
  % references/kleiman-picard-src/kleiman-picard.tex, L3990-L3996).
  % SOURCE QUOTE:
  % "\\label{rmk:Jac}
  %  Assume \(X/S\) is projective and smooth, its geometric fibers are
  % connected curves of genus \(g>0\), and \(S\) is Noetherian.  Set
  % \(J:=\IPicz_{X/S}\); it exists and is a projective Abelian \(S\)-scheme by
  % Exercise~\ref{ex:jac}."
  \textit{Source: \cite{milne-av}, ``Abelian Varieties'' (course notes, 2008),
  \S I.1, p.~8 (definition of abelian variety); \cite{kleiman-picard}, ``The
  Picard scheme'', \S 5, Rmk.~5.26 (the Jacobian
  \(J := \Pic^0_{X/S}\) of a smooth proper family of curves is a
  projective abelian \(S\)-scheme).}
  Let \(C/k\) be a smooth proper geometrically integral curve of
  genus \(g\) over a field \(k\). The identity component
  \(\Pic^0_{C/k}\) (\cref{def:pic_zero_subscheme}) is an
  \emph{abelian variety over \(k\)} of dimension \(g\): the underlying
  \(k\)-scheme is smooth (\cref{thm:pic0_smooth}), proper
  (\cref{thm:pic0_proper}), and geometrically irreducible
  (\cref{thm:pic0_geom_irred}), and it carries a \(k\)-group-scheme
  structure inherited from \(\Pic_{C/k}\) via the inclusion
  \(\Pic^0_{C/k} \hookrightarrow \Pic_{C/k}\)
  (\cref{thm:identity_component_is_subgroup_homomorphism}). In
  particular, \(\Pic^0_{C/k}\) is automatically commutative (Milne~\S
  I.1, Cor.~1.4).
\end{theorem}

\begin{proof}
  \uses{thm:pic0_smooth,
        thm:pic0_proper,
        thm:pic0_geom_irred,
        thm:pic0_group_structure,
        thm:identity_component_is_subgroup_homomorphism,
        thm:identity_component_finite_type_geom_irreducible}
  We assemble Milne's defining axioms of an abelian variety from the
  four conjuncts established earlier in this chapter and in
  \cref{chap:Picard_IdentityComponent}.

  \emph{(i) Group-variety structure.} By
  \cref{thm:identity_component_is_subgroup_homomorphism} applied to
  \(G = \Pic_{C/k}\) (a \(k\)-group scheme locally of finite type by
  \cref{thm:fga_pic_representability}), the identity component
  \(\Pic^0_{C/k}\) inherits a \(k\)-group-scheme structure from
  \(\Pic_{C/k}\) compatible with the clopen inclusion.

  \emph{(ii) Smoothness.} By \cref{thm:pic0_smooth}, the structural
  morphism \(\Pic^0_{C/k} \to \Spec k\) is smooth (of relative
  dimension \(g\)). In particular \(\Pic^0_{C/k}\) is a smooth
  \(k\)-variety: it is reduced (Milne's standing convention bundles
  reducedness into ``variety''), geometrically reduced (by
  smoothness over \(k\)), and locally of finite type.

  \emph{(iii) Properness.} By \cref{thm:pic0_proper}, the structural
  morphism \(\Pic^0_{C/k} \to \Spec k\) is proper (in fact
  projective). This is Milne's ``complete'' axiom.

  \emph{(iv) Connectedness (geometric irreducibility).} By
  \cref{thm:pic0_geom_irred}, the underlying scheme
  \(\Pic^0_{C/k}\) is geometrically irreducible; in particular it is
  geometrically connected, hence connected.

  Combining (i)--(iv): \(\Pic^0_{C/k}\) is a smooth, proper,
  geometrically irreducible \(k\)-scheme equipped with a
  \(k\)-group-scheme structure. By Milne~\S I.1, p.~8, ``a complete
  connected group variety is called an abelian variety''. Thus
  \(\Pic^0_{C/k}\) is an abelian variety over \(k\); commutativity is
  automatic (Milne~\S I.1, Cor.~1.4). Its dimension is the relative
  dimension of the smooth morphism \(\Pic^0_{C/k} \to \Spec k\),
  which equals \(g(C) = \dim_k H^1(C, \mathcal{O}_C)\) by
  \cref{thm:pic0_tangent_space_iso} and \cref{thm:pic0_smooth}.

  This is the content of Kleiman~\S 5 Rmk.~5.26: for \(X/S\)
  projective and smooth with connected curve fibres of genus
  \(g > 0\) (specialise to \(S = \Spec k\), \(X = C\)), the Jacobian
  \(J := \Pic^0_{X/S}\) exists and is a projective abelian
  \(S\)-scheme. The genus-zero case is degenerate
  (\(\Pic^0_{C/k} = 0\), Milne~\S III.1, Rmk.~1.4(e)) but is
  consistent with the assembly: a point scheme \(\Spec k\) is
  trivially an abelian variety of dimension \(0\).
\end{proof}

\section{Lean encoding}
\label{sec:pic0_isAbelianVariety_lean_encoding}

The Lean target file is
\texttt{AlgebraicJacobian/Picard/Pic0AbelianVariety.lean}, sibling to
\texttt{AlgebraicJacobian/Picard/IdentityComponent.lean}
(\cref{chap:Picard_IdentityComponent}). The five theorems of this
chapter correspond to Lean targets as follows.

\begin{itemize}
  \item \texttt{AlgebraicGeometry.Scheme.Pic0.tangentSpaceIso}
    (\cref{thm:pic0_tangent_space_iso}) packages the canonical
    isomorphism
    \(\tu T_0 \Pic^0_{C/k} \cong H^1(C, \mathcal{O}_C)\), the
    curve-specific content of Kleiman~\S 5 Thm.~5.11. In Lean
    this is an isomorphism of \(k\)-modules between the Zariski
    tangent space of \(\Pic^0_{C/k}\) at the identity \(0\) and the
    first sheaf-cohomology \(k\)-module \(H^1(C, \mathcal{O}_C)\).
  \item \texttt{AlgebraicGeometry.Scheme.Pic0.smooth}
    (\cref{thm:pic0_smooth}) packages the smoothness of
    \(\Pic^0_{C/k}\) as a \texttt{[Smooth]} instance on the structural
    morphism, of relative dimension \(g\). Kleiman~\S 5
    Cor.~5.13 + Cor.~5.14.
  \item \texttt{AlgebraicGeometry.Scheme.Pic0.proper}
    (\cref{thm:pic0_proper}) packages the properness of
    \(\Pic^0_{C/k}\) as a \texttt{[IsProper]} instance on the
    structural morphism. Kleiman~\S 5 Thm.~5.4 (projectivity
    upgrade for geometrically normal \(C/k\)).
  \item \texttt{AlgebraicGeometry.Scheme.Pic0.geometricallyIrreducible}
    (\cref{thm:pic0_geom_irred}) packages the geometric
    irreducibility of \(\Pic^0_{C/k}\) as a
    \texttt{[GeometricallyIrreducible]} instance on the structural
    morphism. Kleiman~\S 5 Prp.~5.3 (specialisation of the
    abstract substrate
    \texttt{IdentityComponent.isFiniteTypeGeometricallyIrreducible}
    of \cref{chap:Picard_IdentityComponent} to \(G = \Pic_{C/k}\)).
  \item \texttt{AlgebraicGeometry.Scheme.Pic0.isAbelianVariety}
    (\cref{thm:pic0_isAbelianVariety}) packages the assembled
    abelian-variety structure: the conjunction of \texttt{[Smooth]},
    \texttt{[IsProper]}, \texttt{[GeometricallyIrreducible]} on the
    structural morphism together with a \(k\)-group-scheme structure
    (\texttt{Nonempty (GrpObj (Pic0 C))}) inherited via the clopen
    inclusion of \cref{thm:identity_component_is_subgroup_homomorphism}.
    Milne~\S I.1, p.~8 (defining axioms) + Kleiman~\S 5
    Rmk.~5.26.
\end{itemize}

The five Lean declarations
\texttt{tangentSpaceIso}, \texttt{smooth}, \texttt{proper},
\texttt{geometricallyIrreducible}, and \texttt{isAbelianVariety} are
new project material: they are not currently in Mathlib, and the Lean
skeleton is owed in a follow-up iteration. The downstream
\texttt{isAbelianVariety} consumes the four conjunct theorems together
with the abstract identity-component substrate of
\cref{chap:Picard_IdentityComponent}, and is the load-bearing A.3.vii
gate of Route~A for \texttt{thm:nonempty\_jacobianWitness}.

\section{Out of scope}
\label{sec:pic0_isAbelianVariety_out_of_scope}

This chapter packages \emph{only} the five A.3.iii--vii steps. The
following downstream constructions live in separate sibling chapters
and are explicitly out of scope here:

\begin{itemize}
  \item \textbf{Albanese universal property of \(\Pic^0_{C/k}\)}
    (A.4.d.ii) --- the universal property of \(f_P : C \to \Pic^0_{C/k}\)
    is the content of \cref{chap:Albanese_AlbaneseUP}, gated on this
    chapter via \cref{thm:pic0_isAbelianVariety}.
  \item \textbf{Autoduality \(\Pic^0_{C/k} \cong (\Pic^0_{C/k})^\vee\)}
    (Milne~\S III.6, Thm.~6.6; Kleiman~\S 5 Rmk.~5.26, Abel map) ---
    not needed for the Jacobian's existence as an abelian variety,
    hence out of scope.
  \item \textbf{Torsion component \(\Pic^\tau_{C/k}\)}
    (Kleiman~\S 6) --- for a smooth proper curve
    \(\Pic^\tau_{C/k} = \Pic^0_{C/k}\), so the finer torsion-component
    finiteness statements of Kleiman~\S 6 are not load-bearing on
    Route~A.
  \item \textbf{Genus-zero degenerate case.} When \(g(C) = 0\),
    \(\Pic^0_{C/k}\) is the trivial group scheme \(\Spec k\)
    (Milne~\S III.1, Rmk.~1.4(e)); this is consistent with the
    abelian-variety assembly (a point is trivially an abelian variety
    of dimension \(0\)) and is handled uniformly here, but is not
    separately stated.
\end{itemize}

\section{Internal-consistency check}
\label{sec:pic0_isAbelianVariety_consistency_check}

The five declaration blocks of this chapter form a connected
\texttt{\textbackslash uses}-chain rooted in
\cref{chap:Picard_IdentityComponent},
\cref{chap:Picard_FGAPicRepresentability}, and \cref{chap:Genus}:

\begin{itemize}
  \item \cref{thm:pic0_tangent_space_iso} uses
    \cref{def:pic_zero_subscheme},
    \cref{thm:identity_component_open_subgroup},
    \cref{def:genus} (statement); proof additionally uses
    \cref{thm:fga_pic_representability} and
    \cref{thm:identity_component_base_change_commutes}.
  \item \cref{thm:pic0_smooth} uses
    \cref{thm:pic0_tangent_space_iso},
    \cref{def:pic_zero_subscheme}, and
    \cref{thm:identity_component_is_subgroup_homomorphism} (statement);
    proof additionally uses
    \cref{thm:identity_component_base_change_commutes}.
  \item \cref{thm:pic0_proper} uses
    \cref{def:pic_zero_subscheme} and
    \cref{thm:identity_component_finite_type_geom_irreducible}
    (statement); proof additionally uses
    \cref{thm:fga_pic_representability} and
    \cref{thm:identity_component_base_change_commutes}.
  \item \cref{thm:pic0_geom_irred} uses
    \cref{def:pic_zero_subscheme} and
    \cref{thm:identity_component_finite_type_geom_irreducible}
    (statement); proof additionally uses
    \cref{thm:identity_component_base_change_commutes} and
    \cref{thm:fga_pic_representability}.
  \item \cref{thm:pic0_isAbelianVariety} uses
    \cref{thm:pic0_smooth},
    \cref{thm:pic0_proper},
    \cref{thm:pic0_geom_irred}, and
    \cref{thm:identity_component_is_subgroup_homomorphism}
    (statement); proof additionally uses
    \cref{thm:identity_component_finite_type_geom_irreducible}.
\end{itemize}

The cross-chapter labels
\texttt{def:pic\_zero\_subscheme},
\texttt{thm:identity\_component\_open\_subgroup},
\texttt{thm:identity\_component\_is\_subgroup\_homomorphism},
\texttt{thm:identity\_component\_finite\_type\_geom\_irreducible},
\texttt{thm:identity\_component\_base\_change\_commutes},
\texttt{thm:fga\_pic\_representability}, and
\texttt{def:genus} all exist in sibling chapters
(\texttt{Picard\_IdentityComponent.tex},
\texttt{Picard\_FGAPicRepresentability.tex}, and \texttt{Genus.tex} at
writing time).
